\textbf{Phần 1. Trắc nghiệm lựa chọn}
\setcounterpageref{ex}{0}

\begin{ex} %[2H5V25]
Cho $4$ điểm phân biệt cùng nằm trên một đường thẳng. Hỏi có tất cả bao nhiêu đoạn thẳng được tạo thành từ $4$ điểm đó?
\choice
{$4$ đoạn thẳng}
{$5$ đoạn thẳng}
{\True $6$ đoạn thẳng}
{$8$ đoạn thẳng}
\loigiai{
Gọi 4 điểm là $A, B, C, D$. Các đoạn thẳng là: $AB, AC, AD, BC, BD, CD$. Có tất cả $6$ đoạn thẳng.
\textbf{Đáp án của câu hỏi là C}
}
\end{ex}

\begin{ex} %[2H5V25]
Một đoạn dây đồng uốn thành hình tam giác có độ dài ba cạnh lần lượt là $12\text{ cm}, 15\text{ cm}, 18\text{ cm}$. Nếu duỗi thẳng đoạn dây đó ra thì chiều dài của nó là
\choice
{\True $45\text{ cm}$}
{$43\text{ cm}$}
{$44\text{ cm}$}
{$46\text{ cm}$}
\loigiai{
Chiều dài đoạn dây bằng chu vi tam giác: $12 + 15 + 18 = 45\text{ (cm)}$.
\textbf{Đáp án của câu hỏi là A}
}
\end{ex}

\begin{ex} %[2H5V25]
Trên đoạn thẳng $MN$ dài $24\text{ cm}$, người ta lấy điểm $P$ sao cho $MP = 9\text{ cm}$. Độ dài đoạn thẳng $PN$ là
\choice
{$13\text{ cm}$}
{$14\text{ cm}$}
{$16\text{ cm}$}
{\True $15\text{ cm}$}
\loigiai{
Độ dài $PN = MN - MP = 24 - 9 = 15\text{ (cm)}$.
\textbf{Đáp án của câu hỏi là D}
}
\end{ex}

\begin{ex} %[2H5V25]
Cho ba điểm $A, B, C$ thẳng hàng theo thứ tự đó. Biết $AB = 25\text{ cm}$, đoạn $BC$ ngắn hơn đoạn $AB$ là $8\text{ cm}$. Độ dài đoạn thẳng $AC$ là
\choice
{$40\text{ cm}$}
{\True $42\text{ cm}$}
{$41\text{ cm}$}
{$43\text{ cm}$}
\loigiai{
$BC = 25 - 8 = 17\text{ (cm)}$. Độ dài $AC = AB + BC = 25 + 17 = 42\text{ (cm)}$.
\textbf{Đáp án của câu hỏi là B}
}
\end{ex}

\begin{ex} %[2H5V25]
Hình nào dưới đây KHÔNG phải là đường thẳng?
\choice
{Đường kéo dài vô tận về hai phía}
{Đường thẳng đi qua hai điểm}
{\True Đoạn nối hai điểm bị giới hạn ở hai đầu}
{Đường nét thẳng không giới hạn}
\loigiai{
Đoạn thẳng bị giới hạn ở hai đầu, còn đường thẳng không bị giới hạn.
\textbf{Đáp án của câu hỏi là C}
}
\end{ex}

\begin{ex} %[2H5V25]
Cho $3$ điểm không thẳng hàng. Nối từng cặp điểm ta được hình gì?
\choice
{\True Hình tam giác}
{Hình tứ giác}
{Đường thẳng}
{Đoạn thẳng dài}
\loigiai{
Nối 3 điểm không thẳng hàng tạo thành 3 cạnh của một hình tam giác.
\textbf{Đáp án của câu hỏi là A}
}
\end{ex}

\begin{ex} %[2H5V25]
Một con kiến bò từ điểm $A$ đến điểm $B$ theo một đường cong dài $28\text{ cm}$. Một chú bọ rùa bò theo đường thẳng từ $A$ đến $B$ hết $19\text{ cm}$. Đường đi của kiến dài hơn của bọ rùa bao nhiêu xăng-ti-mét?
\choice
{$8\text{ cm}$}
{$10\text{ cm}$}
{$7\text{ cm}$}
{\True $9\text{ cm}$}
\loigiai{
Đường đi của kiến dài hơn: $28 - 19 = 9\text{ (cm)}$.
\textbf{Đáp án của câu hỏi là D}
}
\end{ex}

\begin{ex} %[2H5V25]
Có $5$ điểm phân biệt trên một đường tròn. Cứ qua hai điểm nối được một đoạn thẳng. Có tất cả bao nhiêu đoạn thẳng?
\choice
{$8$ đoạn thẳng}
{\True $10$ đoạn thẳng}
{$12$ đoạn thẳng}
{$15$ đoạn thẳng}
\loigiai{
Số đoạn thẳng: $5 \times 4 : 2 = 10$ đoạn thẳng.
\textbf{Đáp án của câu hỏi là B}
}
\end{ex}

\begin{ex} %[2H5V25]
Đoạn thẳng $CD$ dài $1\text{ dm } 6\text{ cm}$. Cắt bớt đi $7\text{ cm}$ thì đoạn thẳng còn lại dài
\choice
{$8\text{ cm}$}
{$10\text{ cm}$}
{\True $9\text{ cm}$}
{$11\text{ cm}$}
\loigiai{
Đổi $1\text{ dm } 6\text{ cm} = 16\text{ cm}$. Còn lại: $16 - 7 = 9\text{ (cm)}$.
\textbf{Đáp án của câu hỏi là C}
}
\end{ex}

\begin{ex} %[2H5V25]
Điểm $I$ là điểm nằm giữa hai điểm $A$ và $B$. Biết $AI = 12\text{ cm}$ và $IB = 12\text{ cm}$. Khẳng định nào sau đây là đúng?
\choice
{\True $I$ là điểm chính giữa của đoạn thẳng $AB$}
{$I$ không thuộc đoạn thẳng $AB$}
{$AB$ dài $12\text{ cm}$}
{$AI$ dài hơn $IB$}
\loigiai{
Vì $I$ nằm giữa và $AI = IB = 12\text{ cm}$ nên $I$ là điểm chính giữa đoạn thẳng $AB$.
\textbf{Đáp án của câu hỏi là A}
}
\end{ex}

\begin{ex} %[2H5V25]
Một hàng rào có $6$ chiếc cọc cắm thẳng hàng cách đều nhau. Khoảng cách giữa hai cọc liên tiếp là $3\text{ m}$. Chiều dài hàng rào từ cọc đầu tiên đến cọc cuối cùng là
\choice
{$18\text{ m}$}
{$12\text{ m}$}
{$16\text{ m}$}
{\True $15\text{ m}$}
\loigiai{
Có $6 - 1 = 5$ khoảng cách. Chiều dài hàng rào: $5 \times 3 = 15\text{ (m)}$.
\textbf{Đáp án của câu hỏi là D}
}
\end{ex}

\begin{ex} %[2H5V25]
Cho ba đoạn thẳng có độ dài lần lượt là $14\text{ cm}, 18\text{ cm}, 22\text{ cm}$. Tổng độ dài của ba đoạn thẳng đó là
\choice
{$52\text{ cm}$}
{\True $54\text{ cm}$}
{$53\text{ cm}$}
{$55\text{ cm}$}
\loigiai{
$14 + 18 + 22 = 54\text{ (cm)}$.
\textbf{Đáp án của câu hỏi là B}
}
\end{ex}

\textbf{Phần 2. Trắc nghiệm đúng sai}
\setcounterpageref{ex}{0}

\begin{ex} %[2H5V25]
Xét tính đúng sai của các khẳng định hình học về điểm và đường thẳng:
\choiceTF
{\True Đoạn thẳng có hai điểm mút giới hạn ở hai đầu.}
{Qua hai điểm phân biệt có thể vẽ được vô số đường thẳng.}
{\True Ba điểm cùng nằm trên một đường thẳng được gọi là ba điểm thẳng hàng.}
{\True Độ dài đoạn thẳng luôn là một số lớn hơn 0.}
\loigiai{
\begin{itemize}
    \item Ý a) Đúng: Đúng theo định nghĩa đoạn thẳng.
    \item Ý b) Sai: Qua hai điểm phân biệt chỉ vẽ được duy nhất một đường thẳng.
    \item Ý c) Đúng: Đúng theo khái niệm.
    \item Ý d) Đúng: Đúng.
\end{itemize}
\begin{itemize}
    \item a) \textbf{Đúng}.
    \item b) \textbf{Sai}.
    \item c) \textbf{Đúng}.
    \item d) \textbf{Đúng}.
\end{itemize}
}
\end{ex}

\begin{ex} %[2H5V25]
Cho đoạn thẳng AB dài 20 cm. Điểm M nằm trên đoạn thẳng AB sao cho AM dài 8 cm:
\choiceTF
{\True Điểm M nằm giữa hai điểm A và B.}
{\True Đoạn thẳng MB dài 12 cm.}
{Đoạn thẳng AM dài hơn đoạn thẳng MB.}
{Điểm M là điểm chính giữa của đoạn thẳng AB.}
\loigiai{
\begin{itemize}
    \item Ý a) Đúng: Vì M thuộc đoạn thẳng AB.
    \item Ý b) Đúng: $MB = 20 - 8 = 12\text{ cm}$.
    \item Ý c) Sai: $AM = 8\text{ cm} < MB = 12\text{ cm}$. Khẳng định sai.
    \item Ý d) Sai: Vì $AM \ne MB$ nên M không phải là điểm chính giữa.
\end{itemize}
\begin{itemize}
    \item a) \textbf{Đúng}.
    \item b) \textbf{Đúng}.
    \item c) \textbf{Sai}.
    \item d) \textbf{Sai}.
\end{itemize}
}
\end{ex}

\textbf{Phần 3. Trả lời ngắn}
\setcounterpageref{ex}{0}

\begin{ex} %[2H5V25]
Cho 5 điểm cùng nằm trên một đường thẳng. Hỏi có bao nhiêu đoạn thẳng được tạo thành?
\shortans{10}
\loigiai{
\textbf{Đáp số:} $10$\\
Số đoạn thẳng: $5 \times 4 : 2 = 10$ đoạn thẳng.
}
\end{ex}

\begin{ex} %[2H5V25]
Đoạn thẳng AB dài 34 cm. Điểm C nằm giữa A và B sao cho AC dài 18 cm. Tính độ dài đoạn thẳng CB.
\shortans{16}
\loigiai{
\textbf{Đáp số:} $16$\\
$CB = 34 - 18 = 16\text{ cm}$.
}
\end{ex}

\begin{ex} %[2H5V25]
Người ta cắm 5 cây cọc thẳng hàng cách đều nhau 4 m. Hỏi khoảng cách từ cọc đầu tiên đến cọc cuối cùng là bao nhiêu mét?
\shortans{16}
\loigiai{
\textbf{Đáp số:} $16$\\
Có $5 - 1 = 4$ khoảng cách: $4 \times 4 = 16\text{ m}$.
}
\end{ex}

\begin{ex} %[2H5V25]
Một đoạn dây thép dài 60 cm uốn thành một hình tam giác có hai cạnh dài 18 cm và 24 cm. Tính độ dài cạnh còn lại của hình tam giác.
\shortans{18}
\loigiai{
\textbf{Đáp số:} $18$\\
Cạnh còn lại: $60 - 18 - 24 = 18\text{ cm}$.
}
\end{ex}

\begin{ex} %[2H5V25]
Ba điểm A, B, C thẳng hàng theo thứ tự đó. Biết AB dài 15 cm, AC dài 32 cm. Tính độ dài đoạn thẳng BC.
\shortans{17}
\loigiai{
\textbf{Đáp số:} $17$\\
$BC = 32 - 15 = 17\text{ cm}$.
}
\end{ex}

\begin{ex} %[2H5V25]
Đoạn thẳng MN dài 45 cm, đoạn thẳng PQ ngắn hơn đoạn thẳng MN là 18 cm. Tính độ dài đoạn thẳng PQ.
\shortans{27}
\loigiai{
\textbf{Đáp số:} $27$\\
$45 - 18 = 27\text{ cm}$.
}
\end{ex}
